
\documentclass{article}
\usepackage{colortbl}
\usepackage{makecell}
\usepackage{multirow}
\usepackage{supertabular}

\begin{document}

\newcounter{utterance}

\centering \large Interaction Transcript for game `referencegame', experiment `number\_grids', episode 5 with reflexion{-}Qwen3.5{-}9B.
\vspace{24pt}

{ \footnotesize  \setcounter{utterance}{1}
\setlength{\tabcolsep}{0pt}
\begin{supertabular}{c@{$\;$}|p{.15\linewidth}@{}p{.15\linewidth}p{.15\linewidth}p{.15\linewidth}p{.15\linewidth}p{.15\linewidth}}
    \# & $\;$A & \multicolumn{4}{c}{Game Master} & $\;\:$B\\
    \hline

    \theutterance \stepcounter{utterance}  
    & & \multicolumn{4}{p{0.6\linewidth}}{
        \cellcolor[rgb]{0.9,0.9,0.9}{
            \makecell[{{p{\linewidth}}}]{
                \texttt{\tiny{[P1$\langle$GM]}}
                \texttt{You are given three grids, where each of them is 5 by 5 in size.} \\
\texttt{Grids have empty cells marked with "▢" and filled cells marked with "X".} \\
\texttt{Your task is to generate a referring expression that best describes the target grid while distinguishing it from the two other distractor grids.} \\
\texttt{The first grid is the target grid, and the following two grids are the distractors.} \\
\\ 
\texttt{Target grid:} \\
\\ 
\texttt{□ X X X X} \\
\texttt{□ X □ □ X} \\
\texttt{□ X X X X} \\
\texttt{□ X □ □ X} \\
\texttt{□ X X X X} \\
\\ 
\texttt{Distractor grid 1:} \\
\\ 
\texttt{□ X X X X} \\
\texttt{□ X □ □ X} \\
\texttt{□ X X X X} \\
\texttt{□ □ □ □ □} \\
\texttt{□ X X X X} \\
\\ 
\texttt{Distractor grid 2:} \\
\\ 
\texttt{□ X X X X} \\
\texttt{□ X □ □ X} \\
\texttt{□ X □ □ X} \\
\texttt{□ X □ □ X} \\
\texttt{□ X X X X} \\
\\ 
\texttt{Instruction: Describe the target grid.} \\
\texttt{Generate the referring expression starting with the tag "Expression: " for the given target grid. Omit any other text.} \\
            }
        }
    }
    & \\ \\

    \theutterance \stepcounter{utterance}  
    & \multicolumn{4}{p{0.6\linewidth}}{
        \cellcolor[rgb]{0.9,0.9,0.9}{
            \makecell[{{p{\linewidth}}}]{
                \texttt{\tiny{[P1$\rangle$GM]}}
                \texttt{Expression: The grid where the fourth row contains only empty cells, while the other rows have filled cells in the second, third, and fifth columns.} \\
            }
        }
    }
    & & \\ \\

    \theutterance \stepcounter{utterance}  
    & & & \multicolumn{2}{p{0.3\linewidth}}{
        \cellcolor[rgb]{0.9,0.9,0.9}{
            \makecell[{{p{\linewidth}}}]{
                \texttt{\tiny{[GM$|$GM]}}
                \texttt{Expression: The grid where the fourth row contains only empty cells, while the other rows have filled cells in the second, third, and fifth columns.} \\
            }
        }
    }
    & & \\ \\

    \theutterance \stepcounter{utterance}  
    & & \multicolumn{4}{p{0.6\linewidth}}{
        \cellcolor[rgb]{0.9,0.9,0.9}{
            \makecell[{{p{\linewidth}}}]{
                \texttt{\tiny{[GM$\rangle$P2]}}
                \texttt{You are given three grids, where each of them is 5 by 5 in size.} \\
\texttt{Grids have empty cells marked with "▢" and filled cells marked with "X".} \\
\texttt{You are also given a referring expression that describes one of the given grids.} \\
\texttt{Your task is to select the grid that matches the given referring expression.} \\
\texttt{Generate only the number (in text) of the grid that matches the given expression by selecting first, second, or third.} \\
\\ 
\texttt{First:} \\
\\ 
\texttt{□ X X X X} \\
\texttt{□ X □ □ X} \\
\texttt{□ X □ □ X} \\
\texttt{□ X □ □ X} \\
\texttt{□ X X X X} \\
\\ 
\texttt{Second:} \\
\\ 
\texttt{□ X X X X} \\
\texttt{□ X □ □ X} \\
\texttt{□ X X X X} \\
\texttt{□ □ □ □ □} \\
\texttt{□ X X X X} \\
\\ 
\texttt{Third:} \\
\\ 
\texttt{□ X X X X} \\
\texttt{□ X □ □ X} \\
\texttt{□ X X X X} \\
\texttt{□ X □ □ X} \\
\texttt{□ X X X X} \\
\\ 
\texttt{Expression: The grid where the fourth row contains only empty cells, while the other rows have filled cells in the second, third, and fifth columns.} \\
\texttt{Question: Which grid does the expression refer to?} \\
\texttt{Start with the tag "Answer: ", followed by your selection. Omit any other text.} \\
            }
        }
    }
    & \\ \\

    \theutterance \stepcounter{utterance}  
    & & & \multicolumn{4}{p{0.6\linewidth}}{
        \cellcolor[rgb]{0.9,0.9,0.9}{
            \makecell[{{p{\linewidth}}}]{
                \texttt{\tiny{[GM$\langle$P2]}}
                \texttt{Answer: Second} \\
            }
        }
    }
     \\ \\

    \theutterance \stepcounter{utterance}  
    & & & \multicolumn{2}{p{0.3\linewidth}}{
        \cellcolor[rgb]{0.9,0.9,0.9}{
            \makecell[{{p{\linewidth}}}]{
                \texttt{\tiny{[GM$|$GM]}}
                \texttt{second} \\
            }
        }
    }
    & & \\ \\

    \theutterance \stepcounter{utterance}  
    & \multicolumn{4}{p{0.6\linewidth}}{
        \cellcolor[rgb]{0.95,0.95,0.95}{
            \makecell[{{p{\linewidth}}}]{
                \texttt{\tiny{[P2$\rightarrow$P1]}}
                \texttt{Post{-}episode Reflexion (not sent to GM):} \\
\texttt{New plan: My previous description was ambiguous because it only specified the fourth row and implied the rest of the grid's structure without explicitly defining the pattern for the other rows, which could match the distractors. Specifically, Distractor 1 has a different fourth row, but Distractor 2 has a different pattern in the third and fourth rows. To ensure uniqueness, I must explicitly describe the alternating pattern of the filled cells in the odd{-}numbered rows (rows 1, 3, 5) versus the even{-}numbered rows (rows 2, 4). The correct strategy is to state that rows 1, 3, and 5 are fully filled except for the first and fourth columns, while rows 2 and 4 have empty cells in the third and fourth columns. I will generate an expression that clearly distinguishes the target's specific row{-}by{-}row configuration from both distractors by mentioning the exact column indices of the filled cells for the repeating pattern.} \\
            }
        }
    }
    & & \\ \\

    \theutterance \stepcounter{utterance}  
    & & & \multicolumn{4}{p{0.6\linewidth}}{
        \cellcolor[rgb]{0.95,0.95,0.95}{
            \makecell[{{p{\linewidth}}}]{
                \texttt{\tiny{[P1$\rightarrow$P2]}}
                \texttt{Post{-}episode Reflexion (not sent to GM):} \\
\texttt{New plan: I incorrectly identified the "Second" grid because I misread the fourth row; it contains an empty cell in the first column (□) and filled cells elsewhere, violating the condition that the fourth row must contain *only* empty cells. The "Third" grid is the correct match because its fourth row is entirely empty (□ □ □ □ □), while the other rows have filled cells in the second, third, and fifth columns. Next time, I will carefully verify the entire fourth row of each grid to ensure it is completely empty before confirming the selection.} \\
            }
        }
    }
     \\ \\

\end{supertabular}
}

\end{document}
